
\documentclass{article}
\usepackage{colortbl}
\usepackage{makecell}
\usepackage{multirow}
\usepackage{supertabular}

\begin{document}

\newcounter{utterance}

\centering \large Interaction Transcript for game `function\_detective', experiment `logic\_passive\_labeled\_pairs', episode 231 with Qwen3.6{-}35B{-}A3B{-}FP8{-}without{-}reasoning.
\vspace{24pt}

{ \footnotesize  \setcounter{utterance}{1}
\setlength{\tabcolsep}{0pt}
\begin{supertabular}{c@{$\;$}|p{.15\linewidth}@{}p{.15\linewidth}p{.15\linewidth}p{.15\linewidth}p{.15\linewidth}p{.15\linewidth}}
    \# & \multicolumn{2}{c}{Player} && \multicolumn{2}{c}{Game Master} \\
    \hline

    \theutterance \stepcounter{utterance}  
    & & & \multicolumn{4}{p{0.6\linewidth}}{
        \cellcolor[rgb]{0.9,0.9,0.9}{
            \makecell[{{p{\linewidth}}}]{
                \texttt{\tiny{[P1$\langle$GM]}}
                \texttt{I have a hidden Python function for you to figure out.} \\
\\ 
\texttt{Your job is to infer what it does and then submit a correct Python implementation.} \\
\\ 
\texttt{Function signature:} \\
\texttt{(a: bool, b: bool) {-}> bool} \\
\\ 
\texttt{You have 15 turns total.} \\
\\ 
\texttt{I reveal candidate input{-}output pairs one at a time and tell you whether each one is correct.} \\
\\ 
\texttt{You must make your best SOLVE attempt no later than the second{-}to{-}last turn.} \\
\texttt{Do not wait until the very end.} \\
\\ 
\texttt{You can use these actions:} \\
\texttt{{-} JSON format described at the end of the prompt} \\
\texttt{{-} SOLVE: ```python ... ```} \\
\\ 
\texttt{Do not add import statements in your guess.} \\
\texttt{Anything you might need is already available, so just write the function body in hypothesis and return the answer.} \\
\\ 
\texttt{Preloaded examples:} \\
\texttt{1. (((False, False), True), True)} \\
\\ 
\texttt{Scoring:} \\
\texttt{{-} correctness is required for any non{-}zero score} \\
\texttt{{-} in interactive modes, solving with fewer information requests is better} \\
\\ 
\texttt{Input formatting reminders:} \\
\texttt{{-} strings must be quoted} \\
\texttt{{-} booleans must be True/False} \\
\texttt{{-} use commas between multiple values} \\
\texttt{{-} do not write any extra action tags in the same message} \\
\\ 
\texttt{Format example:} \\
\texttt{You write:} \\
\texttt{\{'current\_mode': <probing/confirming/don't know; choose whichever best describes your current state in the task>, 'input\_rationale': <brief description of why you chose this hypothesis>, 'hypothesis': <either the exact string "unknown" if you have no guess yet OR your current best guess as runnable Python matching the given signature, with a one{-}line docstring stating what it does. Never write placeholder code, if you cannot commit to a guess, use "unknown". Match the given signature.} \\
\texttt{e.g.} \\
\texttt{def f(x):} \\
\texttt{	"""Returns the absolute value of x."""} \\
\texttt{	return x if x >= 0 else {-}x} \\
\\ 
\texttt{Use "unknown" if you have no guess.>, 'confidence\_score': <score between 1{-}100. 100 means highest confidence in your hypothesis>\}} \\
\\ 
\texttt{Output only one valid json per turn.} \\
            }
        }
    }
     \\ \\

    \theutterance \stepcounter{utterance}  
    & \multicolumn{4}{p{0.6\linewidth}}{
        \cellcolor[rgb]{0.561,0.737,0.561}{
            \makecell[{{p{\linewidth}}}]{
                \texttt{\tiny{[P1$\rangle$GM]}}
                \texttt{\{'current\_mode': 'probing', 'input\_rationale': 'I have one example: f(False, True) returns True. This is consistent with several boolean functions like OR, XOR, A implies B, etc. I need more data points to distinguish between them. I will test (False, False) next.', 'hypothesis': 'unknown', 'confidence\_score': 10\}} \\
            }
        }
    }
    & & \\ \\

    \theutterance \stepcounter{utterance}  
    & & & \multicolumn{4}{p{0.6\linewidth}}{
        \cellcolor[rgb]{0.9,0.9,0.9}{
            \makecell[{{p{\linewidth}}}]{
                \texttt{\tiny{[P1$\langle$GM]}}
                \texttt{Game Violation: ParseError: Response is not valid JSON.} \\
\texttt{Output ONLY:} \\
\texttt{One JSON with confidence\_score, current\_mode, input\_rationale, and hypothesis keys ...} \\
\texttt{OR} \\
\texttt{SOLVE: ```python ... ```} \\
            }
        }
    }
     \\ \\

    \theutterance \stepcounter{utterance}  
    & \multicolumn{4}{p{0.6\linewidth}}{
        \cellcolor[rgb]{0.561,0.737,0.561}{
            \makecell[{{p{\linewidth}}}]{
                \texttt{\tiny{[P1$\rangle$GM]}}
                \texttt{\{} \\
\texttt{  "current\_mode": "probing",} \\
\texttt{  "input\_rationale": "I have one example: f(False, True) returns True. This is consistent with several boolean functions like OR, XOR, A implies B, etc. I need more data points to distinguish between them. I will test (False, False) next.",} \\
\texttt{  "hypothesis": "unknown",} \\
\texttt{  "confidence\_score": 10} \\
\texttt{\}} \\
            }
        }
    }
    & & \\ \\

    \theutterance \stepcounter{utterance}  
    & & & \multicolumn{4}{p{0.6\linewidth}}{
        \cellcolor[rgb]{0.9,0.9,0.9}{
            \makecell[{{p{\linewidth}}}]{
                \texttt{\tiny{[P1$\langle$GM]}}
                \texttt{OUTPUT: (((True, True), False), False)} \\
            }
        }
    }
     \\ \\

    \theutterance \stepcounter{utterance}  
    & \multicolumn{4}{p{0.6\linewidth}}{
        \cellcolor[rgb]{0.561,0.737,0.561}{
            \makecell[{{p{\linewidth}}}]{
                \texttt{\tiny{[P1$\rangle$GM]}}
                \texttt{\{} \\
\texttt{  "current\_mode": "probing",} \\
\texttt{  "input\_rationale": "I have two data points: f(False, True) {-}> True and f(True, True) {-}> False. This eliminates OR (which would be True for both) and AND (which would be False for both). It suggests XOR or NAND or similar. Let's test (False, False) to distinguish further. If f(False, False) is False, it's likely XOR. If True, it's likely NAND.",} \\
\texttt{  "hypothesis": "unknown",} \\
\texttt{  "confidence\_score": 20} \\
\texttt{\}} \\
            }
        }
    }
    & & \\ \\

    \theutterance \stepcounter{utterance}  
    & & & \multicolumn{4}{p{0.6\linewidth}}{
        \cellcolor[rgb]{0.9,0.9,0.9}{
            \makecell[{{p{\linewidth}}}]{
                \texttt{\tiny{[P1$\langle$GM]}}
                \texttt{OUTPUT: (((False, True), False), True)} \\
            }
        }
    }
     \\ \\

    \theutterance \stepcounter{utterance}  
    & \multicolumn{4}{p{0.6\linewidth}}{
        \cellcolor[rgb]{0.561,0.737,0.561}{
            \makecell[{{p{\linewidth}}}]{
                \texttt{\tiny{[P1$\rangle$GM]}}
                \texttt{\{} \\
\texttt{  "current\_mode": "probing",} \\
\texttt{  "input\_rationale": "I have three data points: f(False, True) {-}> True, f(True, True) {-}> False, f(False, False) {-}> True. Let's check the remaining case f(True, False). If it is True, the function is XOR. If it is False, it might be NAND (but NAND(False, True) is True, NAND(True, True) is False, NAND(False, False) is True, so NAND fits so far. Wait, NAND(True, False) is True. So both XOR and NAND fit the first three points. I need to test (True, False) to distinguish between XOR and NAND.",} \\
\texttt{  "hypothesis": "unknown",} \\
\texttt{  "confidence\_score": 30} \\
\texttt{\}} \\
            }
        }
    }
    & & \\ \\

    \theutterance \stepcounter{utterance}  
    & & & \multicolumn{4}{p{0.6\linewidth}}{
        \cellcolor[rgb]{0.9,0.9,0.9}{
            \makecell[{{p{\linewidth}}}]{
                \texttt{\tiny{[P1$\langle$GM]}}
                \texttt{OUTPUT: (((True, False), True), False)} \\
            }
        }
    }
     \\ \\

    \theutterance \stepcounter{utterance}  
    & \multicolumn{4}{p{0.6\linewidth}}{
        \cellcolor[rgb]{0.561,0.737,0.561}{
            \makecell[{{p{\linewidth}}}]{
                \texttt{\tiny{[P1$\rangle$GM]}}
                \texttt{\{} \\
\texttt{  "current\_mode": "confirming",} \\
\texttt{  "input\_rationale": "I have four data points: f(False, True) {-}> True, f(True, True) {-}> False, f(False, False) {-}> True, f(True, False) {-}> False. Let's check XOR and NAND again. XOR: F\^{}T=T, T\^{}T=F, F\^{}F=F, T\^{}F=T. This does not match f(False, False){-}>True or f(True, False){-}>False. NAND: \textasciitilde{}(F\&T)=T, \textasciitilde{}(T\&T)=F, \textasciitilde{}(F\&F)=T, \textasciitilde{}(T\&F)=T. This does not match f(True, False){-}>False. Let's re{-}evaluate. The outputs are: (F,T){-}>T, (T,T){-}>F, (F,F){-}>T, (T,F){-}>F. This looks like `not b`. Let's check: not T = F, not T = F (wait, f(T,T) is F, not T is F. f(F,T) is T, not T is F. Mismatch). Let's check `a`. f(F,T)=T (a=F, mismatch). Let's check `b`. f(F,T)=T (b=T, match). f(T,T)=F (b=T, mismatch). Let's check `a != b` (XOR). f(F,T)=T (match). f(T,T)=F (match} \\
            }
        }
    }
    & & \\ \\

    \theutterance \stepcounter{utterance}  
    & & & \multicolumn{4}{p{0.6\linewidth}}{
        \cellcolor[rgb]{0.9,0.9,0.9}{
            \makecell[{{p{\linewidth}}}]{
                \texttt{\tiny{[P1$\langle$GM]}}
                \texttt{Game Violation: ParseError: Response is not valid JSON.} \\
\texttt{Output ONLY:} \\
\texttt{One JSON with confidence\_score, current\_mode, input\_rationale, and hypothesis keys ...} \\
\texttt{OR} \\
\texttt{SOLVE: ```python ... ```} \\
            }
        }
    }
     \\ \\

    \theutterance \stepcounter{utterance}  
    & \multicolumn{4}{p{0.6\linewidth}}{
        \cellcolor[rgb]{0.561,0.737,0.561}{
            \makecell[{{p{\linewidth}}}]{
                \texttt{\tiny{[P1$\rangle$GM]}}
                \texttt{\{} \\
\texttt{  "current\_mode": "probing",} \\
\texttt{  "input\_rationale": "I have four data points: f(False, True) {-}> True, f(True, True) {-}> False, f(False, False) {-}> True, f(True, False) {-}> False. Let's check XOR and NAND again. XOR: F\^{}T=T, T\^{}T=F, F\^{}F=F, T\^{}F=T. This does not match f(False, False){-}>True or f(True, False){-}>False. NAND: \textasciitilde{}(F\&T)=T, \textasciitilde{}(T\&T)=F, \textasciitilde{}(F\&F)=T, \textasciitilde{}(T\&F)=T. This does not match f(True, False){-}>False. Let's re{-}evaluate. The outputs are: (F,T){-}>T, (T,T){-}>F, (F,F){-}>T, (T,F){-}>F. This looks like `not b`. Let's check: not T = F, not T = F (wait, f(T,T) is F, not T is F. f(F,T) is T, not T is F. Mismatch). Let's check `a`. f(F,T)=T (a=F, mismatch). Let's check `b`. f(F,T)=T (b=T, match). f(T,T)=F (b=T, mismatch). Let's check `a != b` (XOR). f(F,T)=T (match). f(T,T)=F (match} \\
            }
        }
    }
    & & \\ \\

    \theutterance \stepcounter{utterance}  
    & & & \multicolumn{4}{p{0.6\linewidth}}{
        \cellcolor[rgb]{0.9,0.9,0.9}{
            \makecell[{{p{\linewidth}}}]{
                \texttt{\tiny{[P1$\langle$GM]}}
                \texttt{Game Violation: ParseError: Response is not valid JSON.} \\
\texttt{Output ONLY:} \\
\texttt{One JSON with confidence\_score, current\_mode, input\_rationale, and hypothesis keys ...} \\
\texttt{OR} \\
\texttt{SOLVE: ```python ... ```} \\
            }
        }
    }
     \\ \\

    \theutterance \stepcounter{utterance}  
    & \multicolumn{4}{p{0.6\linewidth}}{
        \cellcolor[rgb]{0.561,0.737,0.561}{
            \makecell[{{p{\linewidth}}}]{
                \texttt{\tiny{[P1$\rangle$GM]}}
                \texttt{\{} \\
\texttt{  "current\_mode": "confirming",} \\
\texttt{  "input\_rationale": "I have four data points: f(False, True) {-}> True, f(True, True) {-}> False, f(False, False) {-}> True, f(True, False) {-}> False. Let's check XOR and NAND again. XOR: F\^{}T=T, T\^{}T=F, F\^{}F=F, T\^{}F=T. This does not match f(False, False){-}>True or f(True, False){-}>False. NAND: \textasciitilde{}(F\&T)=T, \textasciitilde{}(T\&T)=F, \textasciitilde{}(F\&F)=T, \textasciitilde{}(T\&F)=T. This does not match f(True, False){-}>False. Let's re{-}evaluate. The outputs are: (F,T){-}>T, (T,T){-}>F, (F,F){-}>T, (T,F){-}>F. This looks like `not b`. Let's check: not T = F, not T = F (wait, f(T,T) is F, not T is F. f(F,T) is T, not T is F. Mismatch). Let's check `a`. f(F,T)=T (a=F, mismatch). Let's check `b`. f(F,T)=T (b=T, match). f(T,T)=F (b=T, mismatch). Let's check `a != b` (XOR). f(F,T)=T (match). f(T,T)=F (match} \\
            }
        }
    }
    & & \\ \\

    \theutterance \stepcounter{utterance}  
    & & & \multicolumn{4}{p{0.6\linewidth}}{
        \cellcolor[rgb]{0.9,0.9,0.9}{
            \makecell[{{p{\linewidth}}}]{
                \texttt{\tiny{[P1$\langle$GM]}}
                \texttt{Game Violation: ParseError: Response is not valid JSON.} \\
\texttt{Output ONLY:} \\
\texttt{One JSON with confidence\_score, current\_mode, input\_rationale, and hypothesis keys ...} \\
\texttt{OR} \\
\texttt{SOLVE: ```python ... ```} \\
            }
        }
    }
     \\ \\

    \theutterance \stepcounter{utterance}  
    & \multicolumn{4}{p{0.6\linewidth}}{
        \cellcolor[rgb]{0.561,0.737,0.561}{
            \makecell[{{p{\linewidth}}}]{
                \texttt{\tiny{[P1$\rangle$GM]}}
                \texttt{\{} \\
\texttt{  "current\_mode": "confirming",} \\
\texttt{  "input\_rationale": "I have four data points: f(False, True) {-}> True, f(True, True) {-}> False, f(False, False) {-}> True, f(True, False) {-}> False. Let's check XOR and NAND again. XOR: F\^{}T=T, T\^{}T=F, F\^{}F=F, T\^{}F=T. This does not match f(False, False){-}>True or f(True, False){-}>False. NAND: \textasciitilde{}(F\&T)=T, \textasciitilde{}(T\&T)=F, \textasciitilde{}(F\&F)=T, \textasciitilde{}(T\&F)=T. This does not match f(True, False){-}>False. Let's re{-}evaluate. The outputs are: (F,T){-}>T, (T,T){-}>F, (F,F){-}>T, (T,F){-}>F. This looks like `not b`. Let's check: not T = F, not T = F (wait, f(T,T) is F, not T is F. f(F,T) is T, not T is F. Mismatch). Let's check `a`. f(F,T)=T (a=F, mismatch). Let's check `b`. f(F,T)=T (b=T, match). f(T,T)=F (b=T, mismatch). Let's check `a != b` (XOR). f(F,T)=T (match). f(T,T)=F (match} \\
            }
        }
    }
    & & \\ \\

    \theutterance \stepcounter{utterance}  
    & & & \multicolumn{4}{p{0.6\linewidth}}{
        \cellcolor[rgb]{0.9,0.9,0.9}{
            \makecell[{{p{\linewidth}}}]{
                \texttt{\tiny{[P1$\langle$GM]}}
                \texttt{Game Violation: ParseError: Response is not valid JSON.} \\
\texttt{Output ONLY:} \\
\texttt{One JSON with confidence\_score, current\_mode, input\_rationale, and hypothesis keys ...} \\
\texttt{OR} \\
\texttt{SOLVE: ```python ... ```} \\
            }
        }
    }
     \\ \\

    \theutterance \stepcounter{utterance}  
    & \multicolumn{4}{p{0.6\linewidth}}{
        \cellcolor[rgb]{0.561,0.737,0.561}{
            \makecell[{{p{\linewidth}}}]{
                \texttt{\tiny{[P1$\rangle$GM]}}
                \texttt{\{} \\
\texttt{  "current\_mode": "confirming",} \\
\texttt{  "input\_rationale": "I have four data points: f(False, True) {-}> True, f(True, True) {-}> False, f(False, False) {-}> True, f(True, False) {-}> False. Let's check XOR and NAND again. XOR: F\^{}T=T, T\^{}T=F, F\^{}F=F, T\^{}F=T. This does not match f(False, False){-}>True or f(True, False){-}>False. NAND: \textasciitilde{}(F\&T)=T, \textasciitilde{}(T\&T)=F, \textasciitilde{}(F\&F)=T, \textasciitilde{}(T\&F)=T. This does not match f(True, False){-}>False. Let's re{-}evaluate. The outputs are: (F,T){-}>T, (T,T){-}>F, (F,F){-}>T, (T,F){-}>F. This looks like `not b`. Let's check: not T = F, not T = F (wait, f(T,T) is F, not T is F. f(F,T) is T, not T is F. Mismatch). Let's check `a`. f(F,T)=T (a=F, mismatch). Let's check `b`. f(F,T)=T (b=T, match). f(T,T)=F (b=T, mismatch). Let's check `a != b` (XOR). f(F,T)=T (match). f(T,T)=F (match} \\
            }
        }
    }
    & & \\ \\

    \theutterance \stepcounter{utterance}  
    & & & \multicolumn{4}{p{0.6\linewidth}}{
        \cellcolor[rgb]{0.9,0.9,0.9}{
            \makecell[{{p{\linewidth}}}]{
                \texttt{\tiny{[P1$\langle$GM]}}
                \texttt{Game Violation: ParseError: Response is not valid JSON.} \\
\texttt{Output ONLY:} \\
\texttt{One JSON with confidence\_score, current\_mode, input\_rationale, and hypothesis keys ...} \\
\texttt{OR} \\
\texttt{SOLVE: ```python ... ```} \\
            }
        }
    }
     \\ \\

    \theutterance \stepcounter{utterance}  
    & \multicolumn{4}{p{0.6\linewidth}}{
        \cellcolor[rgb]{0.561,0.737,0.561}{
            \makecell[{{p{\linewidth}}}]{
                \texttt{\tiny{[P1$\rangle$GM]}}
                \texttt{\{} \\
\texttt{  "current\_mode": "confirming",} \\
\texttt{  "input\_rationale": "I have four data points: f(False, True) {-}> True, f(True, True) {-}> False, f(False, False) {-}> True, f(True, False) {-}> False. Let's check XOR and NAND again. XOR: F\^{}T=T, T\^{}T=F, F\^{}F=F, T\^{}F=T. This does not match f(False, False){-}>True or f(True, False){-}>False. NAND: \textasciitilde{}(F\&T)=T, \textasciitilde{}(T\&T)=F, \textasciitilde{}(F\&F)=T, \textasciitilde{}(T\&F)=T. This does not match f(True, False){-}>False. Let's re{-}evaluate. The outputs are: (F,T){-}>T, (T,T){-}>F, (F,F){-}>T, (T,F){-}>F. This looks like `not b`. Let's check: not T = F, not T = F (wait, f(T,T) is F, not T is F. f(F,T) is T, not T is F. Mismatch). Let's check `a`. f(F,T)=T (a=F, mismatch). Let's check `b`. f(F,T)=T (b=T, match). f(T,T)=F (b=T, mismatch). Let's check `a != b` (XOR). f(F,T)=T (match). f(T,T)=F (match} \\
            }
        }
    }
    & & \\ \\

    \theutterance \stepcounter{utterance}  
    & & & \multicolumn{4}{p{0.6\linewidth}}{
        \cellcolor[rgb]{0.9,0.9,0.9}{
            \makecell[{{p{\linewidth}}}]{
                \texttt{\tiny{[P1$\langle$GM]}}
                \texttt{Game Violation: ParseError: Response is not valid JSON.} \\
\texttt{Output ONLY:} \\
\texttt{One JSON with confidence\_score, current\_mode, input\_rationale, and hypothesis keys ...} \\
\texttt{OR} \\
\texttt{SOLVE: ```python ... ```} \\
            }
        }
    }
     \\ \\

    \theutterance \stepcounter{utterance}  
    & \multicolumn{4}{p{0.6\linewidth}}{
        \cellcolor[rgb]{0.561,0.737,0.561}{
            \makecell[{{p{\linewidth}}}]{
                \texttt{\tiny{[P1$\rangle$GM]}}
                \texttt{\{} \\
\texttt{  "current\_mode": "confirming",} \\
\texttt{  "input\_rationale": "I have four data points: f(False, True) {-}> True, f(True, True) {-}> False, f(False, False) {-}> True, f(True, False) {-}> False. Let's check XOR and NAND again. XOR: F\^{}T=T, T\^{}T=F, F\^{}F=F, T\^{}F=T. This does not match f(False, False){-}>True or f(True, False){-}>False. NAND: \textasciitilde{}(F\&T)=T, \textasciitilde{}(T\&T)=F, \textasciitilde{}(F\&F)=T, \textasciitilde{}(T\&F)=T. This does not match f(True, False){-}>False. Let's re{-}evaluate. The outputs are: (F,T){-}>T, (T,T){-}>F, (F,F){-}>T, (T,F){-}>F. This looks like `not b`. Let's check: not T = F, not T = F (wait, f(T,T) is F, not T is F. f(F,T) is T, not T is F. Mismatch). Let's check `a`. f(F,T)=T (a=F, mismatch). Let's check `b`. f(F,T)=T (b=T, match). f(T,T)=F (b=T, mismatch). Let's check `a != b` (XOR). f(F,T)=T (match). f(T,T)=F (match} \\
            }
        }
    }
    & & \\ \\

    \theutterance \stepcounter{utterance}  
    & & & \multicolumn{4}{p{0.6\linewidth}}{
        \cellcolor[rgb]{0.9,0.9,0.9}{
            \makecell[{{p{\linewidth}}}]{
                \texttt{\tiny{[P1$\langle$GM]}}
                \texttt{Game Violation: ParseError: Response is not valid JSON.} \\
\texttt{Output ONLY:} \\
\texttt{One JSON with confidence\_score, current\_mode, input\_rationale, and hypothesis keys ...} \\
\texttt{OR} \\
\texttt{SOLVE: ```python ... ```} \\
            }
        }
    }
     \\ \\

    \theutterance \stepcounter{utterance}  
    & \multicolumn{4}{p{0.6\linewidth}}{
        \cellcolor[rgb]{0.561,0.737,0.561}{
            \makecell[{{p{\linewidth}}}]{
                \texttt{\tiny{[P1$\rangle$GM]}}
                \texttt{\{} \\
\texttt{  "current\_mode": "confirming",} \\
\texttt{  "input\_rationale": "I have four data points: f(False, True) {-}> True, f(True, True) {-}> False, f(False, False) {-}> True, f(True, False) {-}> False. Let's check XOR and NAND again. XOR: F\^{}T=T, T\^{}T=F, F\^{}F=F, T\^{}F=T. This does not match f(False, False){-}>True or f(True, False){-}>False. NAND: \textasciitilde{}(F\&T)=T, \textasciitilde{}(T\&T)=F, \textasciitilde{}(F\&F)=T, \textasciitilde{}(T\&F)=T. This does not match f(True, False){-}>False. Let's re{-}evaluate. The outputs are: (F,T){-}>T, (T,T){-}>F, (F,F){-}>T, (T,F){-}>F. This looks like `not b`. Let's check: not T = F, not T = F (wait, f(T,T) is F, not T is F. f(F,T) is T, not T is F. Mismatch). Let's check `a`. f(F,T)=T (a=F, mismatch). Let's check `b`. f(F,T)=T (b=T, match). f(T,T)=F (b=T, mismatch). Let's check `a != b` (XOR). f(F,T)=T (match). f(T,T)=F (match} \\
            }
        }
    }
    & & \\ \\

    \theutterance \stepcounter{utterance}  
    & & & \multicolumn{4}{p{0.6\linewidth}}{
        \cellcolor[rgb]{0.9,0.9,0.9}{
            \makecell[{{p{\linewidth}}}]{
                \texttt{\tiny{[P1$\langle$GM]}}
                \texttt{Game Violation: ParseError: Response is not valid JSON.} \\
\texttt{Output ONLY:} \\
\texttt{One JSON with confidence\_score, current\_mode, input\_rationale, and hypothesis keys ...} \\
\texttt{OR} \\
\texttt{SOLVE: ```python ... ```} \\
            }
        }
    }
     \\ \\

    \theutterance \stepcounter{utterance}  
    & \multicolumn{4}{p{0.6\linewidth}}{
        \cellcolor[rgb]{0.561,0.737,0.561}{
            \makecell[{{p{\linewidth}}}]{
                \texttt{\tiny{[P1$\rangle$GM]}}
                \texttt{\{} \\
\texttt{  "current\_mode": "confirming",} \\
\texttt{  "input\_rationale": "I have four data points: f(False, True) {-}> True, f(True, True) {-}> False, f(False, False) {-}> True, f(True, False) {-}> False. Let's check XOR and NAND again. XOR: F\^{}T=T, T\^{}T=F, F\^{}F=F, T\^{}F=T. This does not match f(False, False){-}>True or f(True, False){-}>False. NAND: \textasciitilde{}(F\&T)=T, \textasciitilde{}(T\&T)=F, \textasciitilde{}(F\&F)=T, \textasciitilde{}(T\&F)=T. This does not match f(True, False){-}>False. Let's re{-}evaluate. The outputs are: (F,T){-}>T, (T,T){-}>F, (F,F){-}>T, (T,F){-}>F. This looks like `not b`. Let's check: not T = F, not T = F (wait, f(T,T) is F, not T is F. f(F,T) is T, not T is F. Mismatch). Let's check `a`. f(F,T)=T (a=F, mismatch). Let's check `b`. f(F,T)=T (b=T, match). f(T,T)=F (b=T, mismatch). Let's check `a != b` (XOR). f(F,T)=T (match). f(T,T)=F (match} \\
            }
        }
    }
    & & \\ \\

    \theutterance \stepcounter{utterance}  
    & & & \multicolumn{4}{p{0.6\linewidth}}{
        \cellcolor[rgb]{0.9,0.9,0.9}{
            \makecell[{{p{\linewidth}}}]{
                \texttt{\tiny{[P1$\langle$GM]}}
                \texttt{Game Violation: ParseError: Response is not valid JSON.} \\
\texttt{Output ONLY:} \\
\texttt{One JSON with confidence\_score, current\_mode, input\_rationale, and hypothesis keys ...} \\
\texttt{OR} \\
\texttt{SOLVE: ```python ... ```} \\
            }
        }
    }
     \\ \\

    \theutterance \stepcounter{utterance}  
    & \multicolumn{4}{p{0.6\linewidth}}{
        \cellcolor[rgb]{0.561,0.737,0.561}{
            \makecell[{{p{\linewidth}}}]{
                \texttt{\tiny{[P1$\rangle$GM]}}
                \texttt{\{} \\
\texttt{  "current\_mode": "confirming",} \\
\texttt{  "input\_rationale": "I have four data points: f(False, True) {-}> True, f(True, True) {-}> False, f(False, False) {-}> True, f(True, False) {-}> False. Let's check XOR and NAND again. XOR: F\^{}T=T, T\^{}T=F, F\^{}F=F, T\^{}F=T. This does not match f(False, False){-}>True or f(True, False){-}>False. NAND: \textasciitilde{}(F\&T)=T, \textasciitilde{}(T\&T)=F, \textasciitilde{}(F\&F)=T, \textasciitilde{}(T\&F)=T. This does not match f(True, False){-}>False. Let's re{-}evaluate. The outputs are: (F,T){-}>T, (T,T){-}>F, (F,F){-}>T, (T,F){-}>F. This looks like `not b`. Let's check: not T = F, not T = F (wait, f(T,T) is F, not T is F. f(F,T) is T, not T is F. Mismatch). Let's check `a`. f(F,T)=T (a=F, mismatch). Let's check `b`. f(F,T)=T (b=T, match). f(T,T)=F (b=T, mismatch). Let's check `a != b` (XOR). f(F,T)=T (match). f(T,T)=F (match} \\
            }
        }
    }
    & & \\ \\

    \theutterance \stepcounter{utterance}  
    & & & \multicolumn{4}{p{0.6\linewidth}}{
        \cellcolor[rgb]{0.9,0.9,0.9}{
            \makecell[{{p{\linewidth}}}]{
                \texttt{\tiny{[P1$\langle$GM]}}
                \texttt{Game Violation: ParseError: Response is not valid JSON.} \\
\texttt{Output ONLY:} \\
\texttt{One JSON with confidence\_score, current\_mode, input\_rationale, and hypothesis keys ...} \\
\texttt{OR} \\
\texttt{SOLVE: ```python ... ```} \\
            }
        }
    }
     \\ \\

    \theutterance \stepcounter{utterance}  
    & \multicolumn{4}{p{0.6\linewidth}}{
        \cellcolor[rgb]{0.561,0.737,0.561}{
            \makecell[{{p{\linewidth}}}]{
                \texttt{\tiny{[P1$\rangle$GM]}}
                \texttt{\{} \\
\texttt{  "current\_mode": "confirming",} \\
\texttt{  "input\_rationale": "I have four data points: f(False, True) {-}> True, f(True, True) {-}> False, f(False, False) {-}> True, f(True, False) {-}> False. Let's check XOR and NAND again. XOR: F\^{}T=T, T\^{}T=F, F\^{}F=F, T\^{}F=T. This does not match f(False, False){-}>True or f(True, False){-}>False. NAND: \textasciitilde{}(F\&T)=T, \textasciitilde{}(T\&T)=F, \textasciitilde{}(F\&F)=T, \textasciitilde{}(T\&F)=T. This does not match f(True, False){-}>False. Let's re{-}evaluate. The outputs are: (F,T){-}>T, (T,T){-}>F, (F,F){-}>T, (T,F){-}>F. This looks like `not b`. Let's check: not T = F, not T = F (wait, f(T,T) is F, not T is F. f(F,T) is T, not T is F. Mismatch). Let's check `a`. f(F,T)=T (a=F, mismatch). Let's check `b`. f(F,T)=T (b=T, match). f(T,T)=F (b=T, mismatch). Let's check `a != b` (XOR). f(F,T)=T (match). f(T,T)=F (match} \\
            }
        }
    }
    & & \\ \\

    \theutterance \stepcounter{utterance}  
    & & & \multicolumn{2}{p{0.3\linewidth}}{
        \cellcolor[rgb]{0.9,0.9,0.9}{
            \makecell[{{p{\linewidth}}}]{
                \texttt{\tiny{[GM$|$GM]}}
                \texttt{Game Verdict: LOSS (turn limit reached)} \\
            }
        }
    }
    & & \\ \\

\end{supertabular}
}

\end{document}
